\section{Related Work}
\label{sec:related_work}

\subsection{Contamination Detection}

Recent work has produced diverse signals for identifying benchmark data in LLM training corpora. Min-K\% \citep{shi_2023_minkprob} and its successor Min-K\%++ \citep{zhang_2024_minkpp} perform membership inference via token-level surprisal statistics, with the latter achieving stronger separation through normalized probability comparisons. ConStat \citep{dekoninck2024constat} measures accuracy gaps between original and paraphrased items to isolate surface memorization from genuine understanding. CDD and TED \citep{dong_2024_cdd} distinguish generalization from memorization through completion-based probes that test whether models can reproduce exact benchmark content. KDS \citep{choi_2025_kds} quantifies leakage via kernel divergence between model outputs on seen and unseen distributions. TS-Guessing \citep{deng_2024_tsguessing} identifies contamination through systematic evaluation of completion accuracy across temporal splits, revealing that many high-scoring models exhibit suspicious temporal patterns. Concurrent with our work, \citet{tao2025detecting} propose an entropy collapse method targeting RL-trained models, exploiting policy collapse across reasoning paths rather than binary self-evaluation confidence. These methods flag contamination at the item or benchmark level but do not diagnose domain-level effects or produce adjusted scores. ContamCal consumes detection signals as inputs (Stage~1) and feeds them into a calibration framework (Stage~2) that estimates decontaminated ability.

\subsection{IRT for LLM Evaluation}

Item Response Theory (IRT) has gained traction as a principled alternative to raw accuracy for comparing LLMs. \citet{lalor_2024_irt_nlp} provide a comprehensive tutorial on IRT in NLP evaluation, covering model selection, parameter estimation, and ability scoring. TinyBenchmarks \citep{polo_2024_tinybenchmarks} exploits IRT-informed item selection to reduce evaluation cost while preserving ranking fidelity across diverse model families. Growing Pains \citep{habba_2025_growing_pains} validates IRT at scale on 400+ models with fixed parameter calibration and identifies practical limitations in low-respondent regimes where item parameters become poorly constrained. HiBayES \citep{luettgau_2025_hibayes} extends IRT with hierarchical Bayesian inference to produce credible intervals over ability estimates, enabling principled model comparison under uncertainty. ATLAS \citep{li_2024_atlas} applies IRT for adaptive testing and selects maximally informative items to reduce per-model evaluation cost. Across this body of work, all items are treated as equally trustworthy: none models contamination as a covariate within the item response function. ContamCal augments IRT with per-domain contamination parameters $\gamma_d$ that separate genuine ability from exposure-driven inflation.

\subsection{Item Compromise in Educational Measurement}

The psychometric literature has long studied item compromise, where test content is disclosed before administration. \citet{eckerly_2017_item_compromise} develops IRT models that incorporate compromise indicators to jointly estimate item parameters and preknowledge effects, showing that ignoring compromise biases all estimated parameters. \citet{sinharay_2017_preknowledge} proposes likelihood ratio and score tests for item preknowledge detection at the individual level, establishing statistical power requirements for reliable identification. DIF cancellation \citep{ercikan_2008_dif_cancellation} demonstrates that item-level bias can aggregate unpredictably across test domains, motivating domain-level rather than item-level analysis. Explanatory IRT \citep{deboeck_2004_explanatory_irt} provides the theoretical foundation for embedding external covariates directly within item response functions. ContamCal adapts this tradition with three modifications for LLM evaluation: continuous exposure scores replace binary leaked/not-leaked indicators, effects are estimated at domain granularity rather than per-item, and counterfactual inference (setting exposure to zero) replaces group comparison, since LLMs lack a clean ``control group'' (Section~\ref{sec:formulation}).

\subsection{Contamination Mitigation and Score Adjustment}

Mitigation strategies attempt to neutralize contamination at test time. \citet{sun2025emperor} systematically evaluate twenty approaches (item rephrasing, option shuffling, dynamic construction, filtering) across ten LLMs and five benchmarks; none reliably eliminates score inflation. This failure motivates post-hoc statistical correction rather than prevention. DCR \citep{xu2025dcr} is the most closely related method: it estimates decontaminated scores by combining detection signals with fuzzy inference rules. However, DCR produces point estimates without uncertainty quantification, does not model domain-level heterogeneity in contamination effects, and lacks IRT grounding that would connect score adjustment to a latent ability scale. No public implementation is available, so we compare qualitatively in Section~\ref{sec:results}. ContamCal addresses all three limitations through Bayesian IRT with per-domain $\gamma_d$ parameters and posterior credible intervals that propagate signal noise into final uncertainty bounds.

\subsection{Calibration Disambiguation}

We use ``calibration'' in the psychometric sense throughout this paper: adjusting benchmark scores to approximate clean (uncontaminated) performance. This differs from probability calibration in machine learning \citep{guo_2017_calibration}, where the goal is aligning predicted confidences with empirical frequencies via temperature scaling or Platt scaling. ContamCal performs score-level adjustment for contamination bias, not recalibration of output probabilities. We adopt this terminology following the educational measurement convention where ``calibration'' refers to placing items and examinees on a common scale after accounting for known distortions. Readers familiar with the ML calibration literature should note that our reliability diagrams (\S\ref{sec:results}) compare calibrated scores against held-out clean accuracy, not predicted probabilities against observed outcomes.